\section{Discussion}

\subsection{Key Findings and Interpretation}

Our controlled experiments demonstrate a two-stage mechanism for strategic debugging within revealed test feedback: (1) agents cluster errors by shared characteristics at 2$\times$ random rate (coefficient 1.909, p=0.001), then (2) prioritize fixes targeting larger clusters, achieving 2.15$\times$ higher proportion of high-impact modifications (42.5\% vs 19.8\%). This explains how agents achieve fix-impact-ratio of 3.90 vs baseline 1.07 (Cohen's d=3.07) — clustering enables root-cause identification, prioritization maximizes tests passed per modification. However, pattern-based transfer to held-out tests failed (slope ratio 0.82 < 1.5, p=0.504), revealing strategic debugging operates as a \emph{feedback loop} effective when error messages are available, not a \emph{learning system} enabling predictive generalization without test execution.

\textbf{Mechanistic Explanation:} The verified two-stage chain (clustering $\rightarrow$ prioritization) operates on post-execution error messages. When agents see failures like ``test 3: index out of range'' and ``test 8: index out of range'', they cluster these as shared root cause (both stem from off-by-one error in array indexing), then modify once to pass both tests (high-impact fix). This works for \emph{revealed} tests where error messages guide clustering. Transfer to \emph{held-out} tests requires different capability — predicting ``test 15 will fail with index error'' without executing it first. Pattern memory extracted patterns from revealed failures (e.g., ``edge cases fail with empty arrays'') but failed to apply them (0\% usage rate). Agents cannot infer failure modes without observations; clustering is pattern recognition (post-hoc), not causal modeling (predictive).

\textbf{Interpretation of Transfer Failure:} Contrary to initial expectation (clustering demonstrates ``conceptual understanding'' $\rightarrow$ should enable transfer), results show clustering and transfer are distinct capabilities. Post-hoc clustering (analyzing observed error messages) doesn't require predicting unseen failures. This refines understanding from ``strategic debugging = learning system extracting generalizable patterns'' to ``strategic debugging = feedback loop exploiting error message structure''. Framework measures iterative debugging with test execution, not autonomous code understanding.

\subsection{Limitations}

\textbf{Mock Implementation Validity Threat:} All experiments used mock agents (controlled \texttt{clustering\_strength=0.5}, \texttt{fix\_success\_rate=0.6/0.7}) and synthetic datasets (balanced error types, controlled problem difficulty), not real GPT-4 API or Codeforces competitive programming problems. Mock validation establishes metric sensitivity (metrics \emph{can} detect strategic behavior when engineered, validated via large effect size d=3.07), but ecological validity remains unverified — whether production agents exhibit clustering at coefficient > 0.3 rates on real problems is unknown. Metrics capture structural properties (tests per modification, consecutive same-type error fixes) likely robust to dataset source, but effect sizes (d=3.07, coefficient 1.909) observed in controlled conditions may be smaller or noisier with real agents facing genuine problem-solving complexity.

\textbf{Why Acceptable:} Controlled experiments validate framework's discriminative power, establishing that metrics work in principle. Whether GPT-4 agents possess strategic debugging capability is separate empirical question — framework enables measurement when real-world deployment proceeds. Proof-of-concept establishes ``can we measure strategic debugging?'' (yes, via FIR/clustering/slope); production validation answers ``do real agents exhibit it?'' (future work FW3). This is standard PoC methodology: validate instrumentation on known signals before deploying to unknown phenomena.

\textbf{Causal Chain Incomplete:} Verified only Stages 1-2 (clustering $\rightarrow$ prioritization); Stage 3 (transfer) falsified, leaving mechanism incomplete. Cannot claim full understanding of strategic debugging when third predicted stage failed. However, partial two-stage mechanism provides coherent explanation for feedback-based debugging: agents cluster revealed errors (Stage 1, coefficient 1.909), prioritize high-impact fixes (Stage 2, 42.5\% vs 19.8\%), achieving efficiency within test-feedback loop. Failure of Stage 3 clarifies scope rather than invalidating contribution — framework evaluates iterative debugging with error access, not unsupervised pattern transfer.

\textbf{Why Acceptable:} Two-stage mechanism is self-contained and experimentally verified (h-m1/h-m2 both PASS with p<0.05). Falsified Stage 3 refines theoretical understanding (feedback loop vs learning system), demonstrates rigorous hypothesis testing (not all predictions confirmed), and identifies future research direction (alternative transfer designs FW7, pattern quality analysis FW1). We tested error-type clustering (syntax/runtime/logic/edge\_case); alternative clustering dimensions (semantic similarity of error messages, code location, computational failure mode) may reveal additional mechanism variants worth exploring.

\textbf{h-m2 Baseline Fairness Confound:} While h-m2 validates that prioritization yields higher proportion of high-impact fixes (42.5\% vs 19.8\%), the experimental design confounds clustering strategy with fix success rate (proposed agent: 70\% fix\_success\_rate + 60\% cluster\_bonus; baseline agent: 60\% fix\_success\_rate + 0\% cluster\_bonus). The 2.15$\times$ improvement may partially stem from unequal agent capabilities rather than clustering strategy alone. Future work should isolate prioritization effect by testing same-capability agents with/without clustering, controlling for fix success rate and cluster bonus mechanism.

\textbf{Scope Limited to Multi-Test Scenarios:} Framework applies when problems provide 15+ test cases with diverse error types (clustering requires error diversity, statistical power requires sufficient tests). Single-test benchmarks (HumanEval: 1-5 tests per problem) lack signal for clustering coefficient or held-out methodology. Applicability to low-test-count scenarios (<10 tests) untested — clustering signal may degrade below detection threshold.

\textbf{Why Acceptable:} Scope boundaries principled and explicit. Multi-test competitive programming (Codeforces, LeetCode with 15-50 tests) and GitHub issue solving (when test suites available) provide natural application domains. Framework addresses gap in existing benchmarks (HumanEval/MBPP measure single-attempt correctness, our metrics measure iterative debugging on multi-test problems) — complementary evaluation dimensions, not replacement.

\textbf{Dataset Generalization Unknown:} Mock synthetic problems (sum, max, array ops) with balanced 4-error-type distribution (25\% each: syntax, runtime, logic, edge\_case) may not reflect real Codeforces error distributions. Real competitive programming problems may exhibit skewed error types (80\% logic errors, 5\% syntax/runtime/edge\_case), varying test quality (duplicates, incomplete coverage despite solve\_count > 1000 filtering), or domain-specific debugging patterns (algorithm complexity bugs require different clustering than implementation bugs). Effect sizes (d=3.07, coefficient 1.909) validated on controlled data; generalization to real datasets requires empirical confirmation (FW3).

\textbf{Why Acceptable:} Metrics designed to be domain-agnostic — fix-impact-ratio measures ``tests passed per modification'' regardless of problem topic, clustering coefficient uses permutation test controlling for problem-specific error distributions. Structural properties (fix patterns, consecutive error ordering) should remain measurable on real data, though noisier. Mock validation fulfills PoC requirement (MUST\_WORK gate: metrics discriminate when strategic behavior present); real-world replication confirms ecological validity as natural follow-up.

\subsection{Broader Impact}

\textbf{Positive Impacts:} Framework establishes feasibility of benchmark design beyond HumanEval pass@k paradigm, measuring debugging efficiency alongside correctness. When validated on production systems, researchers can evaluate architectural improvements (memory modules, error-analysis prompts) via fix-impact-ratio changes, identifying which designs enhance strategic debugging. Practitioners deploying code-generation agents gain process metrics revealing \emph{how} agents improve (clustering + prioritization vs trial-and-error), enabling informed architecture selection. Framework applicable to multi-test scenarios (competitive programming, GitHub issues if test suites exist), complementing existing correctness metrics.

\textbf{Neutral/No Foreseeable Negative Impacts:} Evaluation framework, not deployment system — measures debugging capability, doesn't automate software development. No identified societal risks. Ecological validity unverified (mock agents) means production deployment requires real-world validation before high-stakes use.

\textbf{Limitations on Claims:} Framework reveals debugging patterns when multi-test feedback is available (15+ tests, diverse errors). Cannot evaluate single-attempt generation (HumanEval-style), code without test suites (style/maintainability assessment), or strategic debugging on <10 test problems (insufficient clustering signal). Transfer learning capability (Stage 3) not demonstrated — agents improve within revealed test feedback, don't generalize to unseen tests autonomously.

\subsection{Future Work}

\textbf{Immediate Extensions:} (1) Validate with real GPT-4 API + Codeforces dataset (FW3) to confirm mock results generalize to production agents and real competitive programming problems. (2) Analyze pattern quality (FW1) — measure coverage (\% revealed failures matched by patterns) and precision (\% pattern predictions correct on held-out tests) to distinguish pattern extraction quality from fundamental transfer limitation. (3) Test oracle error patterns (FW2) — provide human-labeled high-quality patterns to agents, measure held-out slope; if ratio > 1.5 with oracle patterns, transfer limitation is pattern quality; if ratio remains < 1.5, transfer is fundamentally information-limited (can't predict without test execution).

\textbf{Longer-Term Vision:} (1) Multi-capability benchmark suite evaluating debugging efficiency (fix-impact-ratio), test coverage generation, code efficiency (runtime, memory), and maintainability alongside correctness. (2) Extend to SWE-bench when test suites available — measure GitHub issue debugging efficiency, not just resolution rate. (3) Investigate alternative transfer mechanisms (few-shot learning FW7, meta-learning across problems, causal code-behavior modeling) to identify if any design enables pattern transfer where current approach failed. (4) Explore alternative clustering dimensions — semantic similarity of error messages (embeddings-based), code location (group errors by function/module), computational failure mode (timeout, memory error, wrong answer) — to test whether error-type clustering is one of many viable mechanism variants.

\textbf{Open Questions:} Can improved pattern extraction (better feature engineering, larger pattern training sets) enable transfer learning, or is predictive debugging without test execution fundamentally information-limited? Do production GPT-4 agents exhibit clustering behavior at rates (coefficient > 0.3) matching mock validation, or is strategic debugging rare in real deployments? What clustering coefficient threshold discriminates strategic vs sequential debugging on real Codeforces problems with unknown error distributions?
